\documentclass[11pt]{article}

\usepackage[margin=1in]{geometry}
\usepackage{amsmath, amssymb}
\usepackage[T1]{fontenc}
\usepackage[utf8]{inputenc}
\usepackage{lmodern}
\usepackage{enumitem}
\usepackage{hyperref}

\setlist[itemize]{topsep=4pt, itemsep=2pt, parsep=0pt}

\title{6.8610 Spring 2026 Lecture 5 Notes: Evaluation}
\author{Junru Ren with help from Codex}
\date{February 19, 2026}

\begin{document}
\maketitle

\section{Evaluation}

\subsection{Motivation}
This lecture asks a foundational question: \textbf{how should NLP systems actually be evaluated?} Up to this point in the course, many models were trained with maximum likelihood or related objectives, but a low training loss or low perplexity does not necessarily imply that a system is useful, truthful, robust, or aligned with the real task of interest.

\begin{itemize}
    \item The quantity we truly care about is often not directly differentiable or even easy to formalize. The lecture explicitly named goals such as \textbf{truth}, \textbf{beauty}, and \textbf{usefulness}.
    \item Intrinsic metrics such as negative log-likelihood and perplexity are valuable, but they can miss severe practical failures.
    \item Modern NLP systems are evaluated in many ways: compelling demos, benchmark accuracy, human judgments, automatic formal metrics, and increasingly \textbf{language models as judges}.
    \item Evaluation is not only about ranking models; it is also about understanding failure modes, designing trustworthy claims, and choosing metrics that actually reflect the scientific or product goal.
\end{itemize}

\subsection{Key Concepts}
\textbf{Definition: Intrinsic Evaluation}

Intrinsic evaluation measures how well a model fits a formal objective, such as test-set likelihood or perplexity, without necessarily checking performance on a downstream task.

\textbf{Definition: Extrinsic Evaluation}

Extrinsic evaluation measures how well a model helps accomplish a real task, such as translation, classification, summarization, question answering, or interactive assistance.

\begin{itemize}
    \item \textbf{Cool demos:} evaluate whether a system can visibly perform a target task at all. Demos can be motivating and meaningful, but they are also easy to overfit.
    \item \textbf{Forced-choice evaluation:} convert a task into multiple-choice, true/false, or other constrained outputs so scoring is straightforward.
    \item \textbf{Human evaluation:} ask people to rate or rank outputs. This remains the gold standard for many generative tasks, but it is expensive and noisy.
    \item \textbf{Formal metrics:} use automatic functions such as perplexity, Bilingual Evaluation Understudy (BLEU), Recall-Oriented Understudy for Gisting Evaluation (ROUGE), or BERTScore, which uses contextual embeddings from Bidirectional Encoder Representations from Transformers (BERT), to compare predictions with references.
    \item \textbf{LM-as-a-judge:} use a stronger language model to score or compare outputs, often with access to the prompt, a reference answer, a checklist, or external tools.
    \item \textbf{Representative data:} the evaluation set must resemble the real deployment problem closely enough to support the claim being made.
    \item \textbf{Inter-annotator disagreement:} humans do not always agree, especially on nuanced judgments involving style, fluency, truthfulness, or preference.
    \item \textbf{Clever Hans effect:} a model may appear to solve the intended task while actually exploiting spurious cues.
    \item \textbf{Goodhart's Law:} once a metric becomes a target, systems may optimize that metric without improving the underlying goal.
    \item \textbf{Pareto frontier thinking:} model quality is multi-dimensional; accuracy alone may hide tradeoffs involving cost, speed, safety, or usability.
\end{itemize}

\subsection{Core Models and Equations}
\textbf{Training by maximum likelihood.} Much of the course so far has used objectives of the form
\[
\hat{\theta} = \arg\max_{\theta} \sum_{(x,y)} \log p_\theta(y \mid x).
\]

\textbf{Test-set perplexity.} On a held-out set of \(n\) examples,
\[
\operatorname{PPL} = \exp\left\{-\frac{1}{n}\sum_{(x,y)} \log p_\theta(y \mid x)\right\}.
\]
Perplexity is the default intrinsic metric for many language models.

\textbf{Why perplexity is not enough.} The lecture considered a model \(LM_3\) that behaves like a strong model \(LM_1\), except that with probability \(0.5\) it emits garbage such as \texttt{asdfjkl;}. On real test data,
\[
p_3(x) = \frac{1}{2}p_1(x)
\]
implies
\[
\frac{1}{n}\sum_x \log p_3(x)
=
\frac{1}{n}\sum_x \log p_1(x) + \log\frac{1}{2}.
\]
Thus the held-out log-likelihood only suffers a constant penalty, even though the model may be catastrophically bad at generation time.

\textbf{BLEU.} BLEU measures \emph{n}-gram precision between a prediction and a reference. For order \(n\),
\[
\operatorname{Precision}_n
=
\frac{\#(\text{\(n\)-grams in reference and prediction})}
{\#(\text{\(n\)-grams in prediction})}.
\]
The standard BLEU score combines multiple \emph{n}-gram precisions with a brevity penalty (BP):
\[
\operatorname{BLEU}
=
\operatorname{BP}\cdot
\exp\left(
\frac{1}{N}\sum_{n=1}^{N}\log \operatorname{Precision}_n
\right).
\]
The lecture emphasized BLEU as a historically important automatic metric that often correlated surprisingly well with human judgments, especially in early machine translation.

\textbf{ROUGE.} ROUGE is recall-oriented. For example, a 2-gram recall version is
\[
\operatorname{ROUGE}\text{-}2
=
\frac{\#(\text{2-grams in reference and prediction})}
{\#(\text{2-grams in reference})}.
\]
ROUGE is especially common in summarization.

\textbf{BERTScore.} BERTScore replaces exact lexical overlap with contextual embedding similarity. If \(r_i\) are reference token embeddings and \(p_j\) are prediction token embeddings, then a recall-style version is
\[
\operatorname{BERTScore}_{\text{recall}}
=
\frac{1}{|r|}
\sum_i \max_j \cos(r_i, p_j).
\]
This allows synonyms or semantically similar paraphrases to receive partial credit even when the exact words differ.

\textbf{Forced-choice accuracy.} When tasks are phrased as discrete options,
\[
\operatorname{Accuracy}
=
\frac{\#(\text{correct choices})}{\#(\text{examples})}.
\]
This metric is simple, but the lecture stressed that \emph{elicitation procedure} matters: asking for the probability of one token, the probability of a whole answer string, or a generated free response can lead to different measured performance.

\subsection{Algorithms / Architectures}
\textbf{1. Demo-based evaluation}
\begin{itemize}
    \item Define a task \(X\) that a reader can understand directly.
    \item Show the system performing the task.
    \item Use the demo to establish that the capability exists at all.
    \item Then ask whether the behavior generalizes beyond the narrow demo setting.
\end{itemize}

\textbf{2. Forced-choice evaluation}
\begin{itemize}
    \item Convert the task into a constrained output space such as multiple choice, yes/no, or true/false.
    \item Decide how the model will be elicited:
    \begin{itemize}
        \item directly compare probabilities of answer options,
        \item compare probabilities of full answer strings,
        \item or force the model to generate and then map the output to a label.
    \end{itemize}
    \item Score with accuracy or a related discrete metric.
    \item Inspect whether success reflects real task understanding or merely answer recognition and test-taking heuristics.
\end{itemize}

\textbf{3. Human evaluation protocol}
\begin{itemize}
    \item Sample representative inputs from the task distribution.
    \item Blind the evaluation whenever possible so raters do not know which system produced which output.
    \item Ask annotators to rank, rate, or compare outputs.
    \item Measure agreement, monitor fatigue, and check annotator quality.
    \item Budget for cost and reproducibility issues.
\end{itemize}

\textbf{4. Formal metric evaluation}
\begin{itemize}
    \item Build a held-out set with reference answers.
    \item Generate predictions from the system.
    \item Compute one or more automatic metrics such as perplexity, BLEU, ROUGE, or BERTScore.
    \item Compare systems under the same preprocessing and tokenization assumptions.
    \item Use multiple metrics when one metric is likely to miss important failure modes.
\end{itemize}

\textbf{5. LM-as-a-judge evaluation}
\begin{itemize}
    \item Select a judge model, ideally one that is stronger or more reliable than the models being evaluated.
    \item Provide the judge with the prompt, candidate output, and optionally a reference answer.
    \item If possible, give the judge a checklist of required properties.
    \item Optionally allow the judge to use external tools or simulated interaction traces.
    \item Aggregate scores or pairwise preferences across examples.
\end{itemize}

\textbf{6. Efficient benchmarking}
\begin{itemize}
    \item Randomly subsample evaluation data while preserving enough statistical power.
    \item Cluster examples and perform stratified subsampling from each cluster.
    \item Fit estimators that infer benchmark-level scores from a smaller subset of evaluated datapoints.
\end{itemize}

\subsection{Important Intuitions from the Professor}
\begin{itemize}
    \item \textbf{Bad evaluation can invalidate otherwise good work.} The lecture gave a concrete warning about student projects that were evaluated by the project members themselves in a non-blind way; this was described plainly as bad science.
    \item \textbf{A compelling demo is not the same as a solved problem.} The professor emphasized several distinct implications that do \emph{not} follow from a strong demo: sustainable progress, product viability, and task completeness.
    \item \textbf{Forced-choice evaluation is often a practical default.} Even though it constrains the task, it remains one of the cleanest ways to get unambiguous scoring on modern benchmarks.
    \item \textbf{Recognition and generation are different skills.} A model may identify the right answer in multiple-choice form without being able to generate the answer reliably on its own.
    \item \textbf{Humans are still the gold standard, but not for free.} Human evaluation is expensive in time and money, affected by fatigue and demographics, and hard to reproduce exactly.
    \item \textbf{LM judges work partly because verification can be easier than generation.} The lecture stressed that this setting is especially natural when a judge has privileged access to a reference solution or a checklist.
    \item \textbf{Nothing beats human grounding for judge examples.} In-context examples can improve an LM judge, but human-generated examples were described as noticeably stronger than model-generated ones.
    \item \textbf{Most evaluation mistakes are really design mistakes.} The important question is not only how to score outputs, but also what data to choose, who should annotate, and what claim the evaluation is supposed to justify.
    \item \textbf{Do not over-index on one metric.} Accuracy, speed, cost, and usability are all relevant. A system can look worse on one benchmark and still represent the better long-term idea.
\end{itemize}

\subsection{Examples}
\begin{itemize}
    \item \textbf{Benchmark examples:} the lecture opened with real contemporary benchmarks such as ButterBench, land-vs.-water coordinate prediction, and VendingBench to show how broad evaluation has become.
    \item \textbf{LM1/LM2/LM3 thought experiment:} a model that occasionally emits garbage may still not look as terrible as expected under held-out log-likelihood, illustrating the limits of perplexity.
    \item \textbf{SHRDLU:} an impressive classic demo system that worked because it operated in a tiny, closed-world domain with hand-built grammars and semantics.
    \item \textbf{GPT-2 and Google Glass:} both were used to separate technical capability from broader questions about product success, social acceptance, and long-term progress.
    \item \textbf{SuperGLUE, a successor to the General Language Understanding Evaluation (GLUE) benchmark, and Humanity's Last Exam:} benchmark examples showing the appeal of forced-choice evaluation and the dangers of overfitting to one public evaluation set.
    \item \textbf{Machine translation:} a reference translation and a predicted translation were used to motivate formal similarity metrics such as BLEU and ROUGE.
    \item \textbf{Representative-data problem:} Europarl parliamentary proceedings versus Haitian crisis messages illustrated that a training or evaluation corpus can be technically rich yet socially unrepresentative of the intended use case.
    \item \textbf{Checklist judging:} an LM judge can be asked whether an answer satisfies requirements such as listing Airbnbs in Singapore, ensuring accommodations fit two people, and respecting a price ceiling.
    \item \textbf{Clever Hans in fake-news detection:} the lecture highlighted work showing that models could predict labels using evidence alone, without looking at the actual claim, revealing spurious shortcuts.
\end{itemize}

\subsection{Common Pitfalls / Limitations}
\begin{itemize}
    \item \textbf{Overfitting to the metric:} optimizing a benchmark score can move the model away from the real task objective.
    \item \textbf{Brittle demos:} a system may succeed only in a tightly constrained environment.
    \item \textbf{Benchmark leakage and overexposure:} public development sets can gradually become part of the optimization loop, even if they are nominally held out.
    \item \textbf{Forced-choice distortion:} not all tasks naturally fit multiple-choice framing, and success may partly reflect answer recognition rather than task mastery.
    \item \textbf{Human-evaluation noise:} raters vary in preferences, energy, demographics, expertise, and even whether they are genuine human annotators.
    \item \textbf{BLEU and ROUGE limitations:} exact overlap metrics miss synonymy, semantic similarity, and the fact that some words matter much more than others.
    \item \textbf{LM judges are not free labels:} they inherit biases, can be costly, require prompt engineering, and may need stronger human grounding than is initially obvious.
    \item \textbf{Weak scientific rigor:} many NLP comparisons do not fully account for randomness over datasets, training runs, or model outputs.
    \item \textbf{Spurious correlations:} models can exploit formatting cues, answer length, position bias, punctuation, or annotation artifacts rather than learning the intended reasoning skill.
\end{itemize}

\subsection{Connections to Other Lectures}
\begin{itemize}
    \item This lecture reframes Lecture 1's emphasis on language modeling: a strong model of \(p_\theta(y \mid x)\) is not automatically a strong system for truth, usefulness, or downstream utility.
    \item It clarifies how to evaluate the models from Lecture 2. Log-linear models and word embeddings can be compared intrinsically, but downstream tasks often demand different metrics.
    \item It connects directly to Lecture 3's seq2seq and attention material, especially through translation and summarization metrics such as BLEU and ROUGE.
    \item It also connects to Lecture 4's BERT and GPT-style models: encoder and decoder systems are now often evaluated with large benchmarks, preference judgments, and LM-as-a-judge pipelines rather than with a single intrinsic metric.
    \item Conceptually, this lecture adds the missing question behind all previous modeling lectures: \emph{what does it mean for the model to be good?}
\end{itemize}

\subsection{Exam-Relevant Summary}
\begin{itemize}
    \item Maximum-likelihood training and held-out perplexity are important, but they are not sufficient for judging usefulness or generation quality.
    \item Perplexity is
    \[
    \operatorname{PPL} = \exp\left\{-\frac{1}{n}\sum_{(x,y)} \log p_\theta(y \mid x)\right\}.
    \]
    \item A model can have tolerable test likelihood while still generating catastrophic outputs at inference time.
    \item Main evaluation families in this lecture: cool demos, forced-choice evals, human evals, formal metrics, and LM judges.
    \item Forced-choice evals are easy to score, but their elicitation procedure matters and they do not fit every task naturally.
    \item Human evaluation is still the gold standard for many generative tasks, but it is expensive, noisy, and hard to reproduce.
    \item BLEU is based on \emph{n}-gram precision; ROUGE is recall-oriented and common in summarization.
    \item BERTScore replaces exact token overlap with contextual embedding similarity.
    \item LM-as-a-judge is useful when verification is easier than generation, especially if the judge has access to references, checklists, or tools.
    \item Good evaluation needs representative inputs and trustworthy baselines. It also requires awareness of benchmark leakage.
    \item Clever Hans effects and Goodhart's Law are central warnings: high scores do not guarantee that the model learned the intended skill.
    \item The right way to compare systems is often multi-dimensional: cost, speed, robustness, and usefulness may matter as much as raw accuracy.
\end{itemize}

\end{document}
